\section{Pathway-group attribution stability under training-only preprocessing}

A fresh six-split audit predicts CDKN2A nonsynonymous mutation status using the native co-mutation GNN. Gene vocabulary, graph adjacency, and missing-age imputation are estimated using training samples only. Each stratified 75/25 split trains for 120 epochs, then permutes each pathway as a group ten times on held-out samples. The target gene is excluded. Pathway-linked clinical covariates are permuted together with the corresponding gene channels; this is an association audit, not a causal intervention.

\begin{center}\begin{tabular}{lrr}\hline Pathway & Mean AUC drop & Positive splits / 6\\\hline

RTK\_RAS & 0.0512 & 6 \\

TP53 & 0.0281 & 3 \\

CELL\_CYCLE & 0.0000 & 0 \\

TGF\_BETA & -0.0055 & 1 \\

SWI\_SNF & 0.0291 & 6 \\

DNA\_REPAIR & 0.0087 & 4 \\

\hline\end{tabular}\end{center}

Held-out AUC ranges from 0.562 to 0.847. The median pairwise pathway-ranking Spearman correlation is 0.371, with range -0.543 to 1.000.

The positive mean RTK/RAS association is a candidate signal, not a validated mechanism. Splits overlap in patients, curated pathway groups overlap in genes, and group permutation can alter relationships with global burden. Cell-cycle importance is zero in these fits; selecting a known pathway does not ensure measurable incremental signal. These results neither establish a clinical predictor nor close the independent-cohort validation gap. The full splits and per-group results are retained in results/pathway\_stability.json.
